% Lamalunga: the theme describes itself. Every number and color on these
% slides is computed when the document is compiled.
\documentclass[11pt]{beamer}

\providecommand\demooptions{}
\edef\next{\noexpand\usetheme[\demooptions]{lamalunga}}\next

\usepackage{booktabs}
\usepackage{listings}
\usepackage{pgfplots}
\pgfplotsset{compat=1.18}

\title{Lamalunga}
\subtitle{A beamer theme of golden proportions and computed colors}
\author{Donato Meoli}
\institute{github.com/dmeoli/lamalunga}
\date{\today}

% a golden swatch with the name and the hexadecimal code of a color
\newcommand\swatch[2]{%
  \begin{tikzpicture}[baseline=(s.base)]
    \fill[#1] (0,0) rectangle (1.618,1);
    \convertcolorspec{named}{#1}{HTML}\swatchhex
    \node[anchor=north west, inner sep=0pt, align=left,
      font=\modularsize{-0.5}] (s) at (0,-0.12)
      {#2\\[-1pt]{\ttfamily\color{lamalunga-muted}\#\swatchhex}};
  \end{tikzpicture}}
\newcommand\harmonyrow[1]{%
  \colorharmony{demo}{#1}%
  \swatch{demo-base}{base}\hfill
  \swatch{demo-complement}{complement}\hfill
  \swatch{demo-split-a}{split}\hfill
  \swatch{demo-triad-a}{triad}\hfill
  \swatch{demo-analog-a}{analog}}

% a length in millimetres
\ExplSyntaxOn
\NewDocumentCommand \inmm { m }
  { \fp_eval:n { round( \dim_to_decimal_in_unit:nn {#1} { 1mm } , 1 ) } \,mm }
\ExplSyntaxOff

\begin{document}

\maketitle

\section{Proportion}

\begin{frame}{The page is a golden rectangle}
  \begin{goldencolumns}
    The page measures \inmm{\paperwidth} by \inmm{\paperheight}, i.e., a
    ratio of
    \(\pgfmathparse{\paperwidth/\paperheight}\pgfmathprintnumber[precision=4]{\pgfmathresult}\),
    and the margins are \fib{8}\,pt and \fib{7}\,pt,
    two consecutive Fibonacci numbers: since their ratio is again
    \(\varphi\), the block of text is a golden rectangle as well.
    \nextcolumn
    \goldenspiral{2.4cm}
  \end{goldencolumns}
\end{frame}

\begin{frame}{Golden columns}
  \begin{goldencolumns}
    \begin{block}{Major}
      The first column takes \(\varphi^{-1} \approx 61.8\%\) of the width,
      the gutter is \fib{8}\,pt.
    \end{block}
    \nextcolumn
    \begin{block}{Minor}
      The second takes \(\varphi^{-2} \approx 38.2\%\).
    \end{block}
  \end{goldencolumns}
  \fibskip{8}
  \begin{goldencolumns}[minor]
    \begin{exampleblock}{Minor first}
      \texttt{[minor]} swaps them.
    \end{exampleblock}
    \nextcolumn
    \begin{alertblock}{Major second}
      The two parts always add up to the text width.
    \end{alertblock}
  \end{goldencolumns}
\end{frame}

\begin{frame}{A modular scale of type}
  Every size is the base times a power of the ratio, here \(\varphi\);
  half steps give the small print.
  \fibskip{7}
  \begin{tabular}{@{}lrl@{}}
    \toprule
    step & size (pt) & used for \\
    \midrule
    \(2\)   & \modularscale{2}    & title, section pages, standout \\
    \(1\)   & \modularscale{1}    & subtitles \\
    \(1/2\) & \modularscale{0.5}  & frame titles \\
    \(0\)   & \modularscale{0}    & text, block titles \\
    \(-1/2\) & \modularscale{-0.5} & foot line, captions, footnotes \\
    \bottomrule
  \end{tabular}
  \fibskip{7}
  {\modularsize{-0.5}\color{lamalunga-muted}
   The half step \(\sqrt{\varphi}\) is the ratio of the sides of the Kepler
   triangle; spacings are Fibonacci numbers \fib{4}, \fib{5}, \fib{6},
   \fib{7}, \fib{8}\,pt.}
\end{frame}

\section{Color}

\begin{frame}{One color, and the others follow}
  \harmonyrow{lamalunga-seed}
  \fibskip{8}
  \begin{goldencolumns}
    The accents are rotations of the hue of the primary on the
    painter's wheel (red, yellow, blue), so violet calls for gold.
    Before they are used for text, they are moved in lightness until
    they reach a contrast of 4.5:1.
    \nextcolumn
    \begin{tabular}{@{}lr@{}}
      \toprule
      on the canvas & ratio \\
      \midrule
      \textcolor{lamalunga-ink}{ink} & \contrastratio{lamalunga-ink}{lamalunga-canvas} \\
      \textcolor{lamalunga-primary}{primary} & \contrastratio{lamalunga-primary}{lamalunga-canvas} \\
      \textcolor{lamalunga-accent}{accent} & \contrastratio{lamalunga-accent}{lamalunga-canvas} \\
      \textcolor{lamalunga-muted}{muted} & \contrastratio{lamalunga-muted}{lamalunga-canvas} \\
      \bottomrule
    \end{tabular}
  \end{goldencolumns}
\end{frame}

\begin{frame}{Any color, by its hexadecimal code}
  \definecolor{demo-red}{HTML}{C0392B}\harmonyrow{demo-red}
  \fibskip{7}
  \definecolor{demo-teal}{HTML}{1F7A7A}\harmonyrow{demo-teal}
  \fibskip{7}
  \definecolor{demo-gold}{HTML}{D4A017}\harmonyrow{demo-gold}
\end{frame}

\begin{frame}{Golden tints}
  Tints dilute a color by \(\varphi^{-k}\) towards the canvas.
  \fibskip{8}
  \foreach \k in {1,...,8}{\swatch{lamalunga-primary-t\k}{t\k}\hfill}\par
  \fibskip{5}
  \foreach \k in {1,...,8}{\swatch{lamalunga-accent-t\k}{t\k}\hfill}
\end{frame}

\section{Composition}

\lamalungaset{guides=both}
\begin{frame}{Thirds and golden sections}
  The dashed lines divide the page in thirds, the solid ones at the golden
  section; \texttt{guides=both} draws them on every frame while one
  composes the slides.
  \focalpoint{se}{\tikz\fill[lamalunga-accent-vivid] circle (3pt);}
\end{frame}
\lamalungaset{guides=none}

\begin{frame}[standout]
  Less, but better.
\end{frame}

\section{Elements}

\begin{frame}{Lists and blocks}
  \begin{goldencolumns}
    \begin{itemize}
      \item squares whose side
      \begin{itemize}
        \item shrinks by \(\varphi\)
        \begin{itemize}
          \item at every level
        \end{itemize}
      \end{itemize}
      \item \alert{alerts} take the accent
    \end{itemize}
    \begin{enumerate}
      \item numbers
      \item in the primary
    \end{enumerate}
    \nextcolumn
    \begin{block}{Block}
      Golden tints.
    \end{block}
    \begin{alertblock}{Alert}
      The accent.
    \end{alertblock}
    \begin{exampleblock}{Example}
      The second accent.
    \end{exampleblock}
  \end{goldencolumns}
\end{frame}

\begin{frame}[fragile]{Mathematics, code and plots}
  \begin{goldencolumns}
    \[
      \varphi = \frac{1 + \sqrt{5}}{2}, \qquad
      F_n = \operatorname{round}\Bigl(\frac{\varphi^n}{\sqrt{5}}\Bigr)
    \]
\begin{lstlisting}[language=Python]
def fib(n):
    # Binet's formula
    phi = (1 + 5 ** 0.5) / 2
    return round(phi ** n / 5 ** 0.5)
\end{lstlisting}
    \nextcolumn
    \begin{tikzpicture}
      \begin{axis}[width=\linewidth, height=0.85\linewidth,
          tick label style={font=\modularsize{-1}}]
        \addplot+[domain=0:3, samples=7] {1.618^x};
        \addplot+[domain=0:3, samples=7] {1.5^x};
      \end{axis}
    \end{tikzpicture}
  \end{goldencolumns}
\end{frame}

\begin{frame}[standout]
  Thank you
\end{frame}

\appendix

\begin{frame}{Backup}
  Frames after \texttt{\textbackslash appendix} are not counted in the
  total.
\end{frame}

\end{document}
